\section{Scope, results, and relation to the existing drafts}
\label{intro:section}

This article continues the research programme in
\path{Analysis/Polylogarithms/docs} of the ProveIt repository.
The inspected snapshot is commit
\texttt{3a6d80ed618d839915deb0d19ab687f45e81d995},
dated 8 October 2026 UTC.  It contains eight articles and thirty-one
reports, covering multiple polylogarithms, cyclotomic values, gamma and
polygamma reductions, Herglotz functions, and generalized Stieltjes
constants \cite{RepoSnapshot}.  The source inventory and exact Git blob
identifiers are recorded in \path{PROVENANCE.json}.

The strongest analytic development here concerns real zeros of the
Stieltjes parameter derivatives.  The existing derivative-tower article
\cite{DraftTower} expresses these functions through Hurwitz zeta jets.
We reorganize that identity as a Laplace transform whose sign-changing
polynomial is independent of the derivative order.  This makes a global
zero bound and a complete large-order description possible.
The strongest algebraic development is a proof of the unrestricted
rational-grid rank formulas that were stated or conjectured from finite
matrices in the two Stieltjes articles \cite{DraftTower,DraftLadder}.

\subsection{The main analytic statements}

Our normalization is
\begin{equation}\label{intro:stieltjes}
 \zeta(s,a)=\frac1{s-1}
       +\sum_{n=0}^{\infty}\frac{(-1)^n}{n!}\gamma_n(a)(s-1)^n,
 \qquad a>0,
\end{equation}
and $\gamma_n^{(k)}$ always means a derivative in $a$.
For fixed Laurent index $n$, define
\[
 P_n(x)=n![z^n]\frac{e^{xz}}{\Gamma(1+z)}.
\]
The results of Sections~\ref{zero:section}--\ref{zero:asymptotic-section}
are the following.

\begin{enumerate}
\item For every $k\geq1$, $\gamma_n^{(k)}$ has at most $n$ positive real
zeros, counted with multiplicity.  This is a global statement on the
whole half-line, not a count inside a numerically chosen window.
\item The polynomial $P_n$ has $n$ distinct real roots.  For every fixed
$n$, the zero bound is attained for all sufficiently large $k$; all
zeros are simple and correspond bijectively to those roots.
\item If $\rho$ is a root of $P_n$ and $t=e^\rho$, its associated zero is
\begin{equation}\label{intro:zero-result}
 a_{\rho,k}=(k+1)e^{-\rho}
   +\frac{e^{-\rho}}2\left(\frac{P_n''(\rho)}{P_n'(\rho)}-1\right)
   -\frac1{\exp(e^\rho)-1}+O_n(k^{-1}).
\end{equation}
We give an explicit $1/(k+1)$ correction, an $O_n(k^{-2})$ remainder
after including it, and a recurrence for the full asymptotic expansion.
\item At $n=1$ there is a unique positive zero for every $k\geq1$, with
the global bracket
\[
 e^{H_{k-1}}<a_k<e^{H_{k-1}}+\tfrac12.
\]
It increases strictly with $k$.  More generally, once all $n$ simple
zeros are present, successive derivative orders strictly interlace.
\end{enumerate}

The elementary mechanism is worth emphasizing.  A reciprocal-gamma
Appell polynomial determines the sign changes of the Laplace kernel.
Rolle's theorem bounds how many zeros its transform can have.  A gamma
probability density then concentrates at $t=(k+1)/a$, yielding all
nearby zeros and their expansions.  Combining the local concentration
argument with the global count excludes zeros on unexamined scales.

\subsection{The formal relation spaces}

For the rational grids, one must first specify which values are unknown
and which are prescribed on the right-hand sides.  Our endpoint symbol
represents the value at $a=1$, not the singular limit at zero.
We prove a weighted distribution theorem for arbitrary rational prime
weights, including zero and resonant weights.  The resulting formal
quotient has dimension $\varphi(q)$ before prescribing the endpoint and
$\varphi(q)-1$ afterwards.  The proof identifies each character block
with its primitive-conductor coordinate and does not divide by an
Euler factor that might vanish.

The same construction works over every commutative coefficient algebra
over a characteristic-zero splitting field. Applied to
$K[\eta]/(\eta^{N+1})$, it proves freeness of the entire finite jet
presentation, with dimension $(N+1)(\varphi(q)-1)$ over $K$ after all
endpoint layers are prescribed. Here $K$ contains the relevant character
values and prime logarithms. This statement concerns the specified
distribution rows alone; additional evaluated identities can reduce a
larger arithmetic presentation.

In particular, the positive Stieltjes derivative grids have formal
residual dimension $\varphi(q)-1$.  For negative zeta jets, after
including the specified reflection identities and endpoint, the
dimension is
\[
 r_{q,k}=\begin{cases}
 \varphi(q)/2-1,& k\ \text{odd},\\
 \varphi(q)/2,& k\ \text{even},
 \end{cases}\qquad(q\geq3).
\]
These statements hold for every positive $k$ and every denominator,
not merely the finite ranges in the drafts.  They are dimensions of
formal presentations.  Injectivity after evaluation as actual complex
numbers is an additional arithmetic question.

We also complete the derivative-tower article's description at trivial
zeros with the exact identity
\[
 \frac{L'(k+1,\chi)}{L(k+1,\chi)}
 +\frac12\overline{\frac{L''(-k,\chi)}{L'(-k,\chi)}}
 =\gamma+\log(2\pi/q)-H_k,
 \qquad\chi(-1)=(-1)^k.
\]
The factor $1/2$ comes from cancelling the gamma pole against the
simple $L$-zero before taking a logarithmic derivative.

\subsection{Classical results, new proofs, and historical priority}

Several central ingredients predate these drafts.  Coffey gives the
general parameter-derivative formula and proves the unique maximum of
$\gamma_1(a)$, the $n=k=1$ case of our zero discussion
\cite{Coffey2009}.  Kubert's universal ordinary distribution gives the
classical rank-$\varphi(q)$ context \cite{Kubert1979}.  The Gaussian
all-even-weight reduction, including the draft's $S_5$ and $S_7$
specializations, appears in Olaikhan's 2022 derivation
\cite{gau-olaikhan-2022}.  Section~\ref{gau:section} gives a complete
analytic reconstruction with an exact generating function and a
convergent resummation, with that attribution retained.

The present continuation establishes the real-zero and asymptotic
theorems just stated, closes the two unrestricted formal-rank gaps by an
explicit weighted proof, and supplies correction proofs and witnesses.
The literature search was targeted to these statements and their
ingredients; it was not an exhaustive priority review.  We therefore
claim the proofs and their application to the repository, without
asserting that every corollary is absent from all earlier literature.
The mathematical arguments were checked in separate derivations during
preparation.  This remains an unrefereed research draft, without a Lean
formalization.

\subsection{Audit policy and source scope}

The audit concerns the retrieved \LaTeX{} and text sources.  All eight
article sources were examined, together with the reports supporting the
claims discussed below.  The remaining reports were inventoried for
context; no assertion of exhaustive verification of every displayed
formula in all thirty-nine documents is made.
Section~\ref{audit:section} provides explicit locations, counterexamples,
or replacement arguments.  The accompanying correction register
distinguishes false formulas from missing hypotheses, incomplete proofs,
and unproved numerical-independence claims.

This distinction matters in this corpus.  An exact rational matrix rank
proves something about that matrix.  It does not prove that every
numerical relation has been included.  A negative integer-relation
search does not establish transcendence or independence.  Nor does a
positive dimension of an ambient motivic space identify the classes of
specific named periods.  The corrected theorems keep these three issues
separate while preserving the valid identities and computations.

\subsection{Notation}

We use $H_k=\sum_{j=1}^k j^{-1}$, $H_0=0$,
$\beta(s)=\sum_{r\geq0}(-1)^r(2r+1)^{-s}$, and $\varphi$ for Euler's
totient.  A prime on $L$ or $\zeta$ means differentiation in the complex
variable $s$ unless a parameter derivative is explicitly indicated.
For $\nu\geq2$, the Clausen convention is
\[
 \operatorname{Cl}_\nu(\theta)=
 \begin{cases}
  \Im\Li_\nu(e^{i\theta}),&\nu\ \text{even},\\
  \Re\Li_\nu(e^{i\theta}),&\nu\ \text{odd}.
 \end{cases}
\]
The Herglotz function in the audit is
$F(x)=\sum_{m\geq1}(\psi(mx)-\log(mx))/m$, unrelated to
$F_{n,k}$ except for the letter used in its source literature.
Each section restates the hypotheses needed for its own formulas.
